\documentclass[10pt,a4paper]{amsart}
\usepackage[margin=19mm]{geometry}
\usepackage{amsmath,amssymb,mathtools}
\usepackage[scr=rsfs,cal=euler]{mathalfa}
\usepackage{xfrac}
\usepackage[unicode,hidelinks,bookmarks=false]{hyperref}
\newcommand{\ie}{i.e.\ }
\newcommand{\cf}{cf.\ }
\newcommand{\resp}{respectively\ }
\newcommand{\Subsection}{\subsection}
\providecommand{\nobreakdash}{\nobreak\hbox{-}}
\setlength{\emergencystretch}{2em}
\multlinegap0pt
\usepackage{amssymb}
\usepackage{float}
\usepackage{tikz}
\newtheorem{theorem}{Theorem}
\title{Singular plane curves: freeness and combinatorics}
\author{Michael Cuntz, Piotr Pokora}
\date{Source: arXiv:2403.13377v5 (CC BY 4.0)}
\begin{document}
\maketitle

\noindent\emph{Source and licence.} Singular plane curves: freeness and combinatorics. Version arXiv:2403.13377v5 (CC BY 4.0), distributed by arXiv under the Creative Commons Attribution 4.0 International licence, http://creativecommons.org/licenses/by/4.0/, as recorded in the arXiv licence metadata for this exact version and shown on the abstract page https://arxiv.org/abs/2403.13377v5. Full licence notice: \texttt{licenses/arxiv-2403.13377-CC-BY-4.0.txt}.\par\smallskip

\noindent\textit{Editorial scope.} Counted unit: the article's theorem pair (source0.tex lines 359--364). The article's complete original proof of it is delivered with it (source0.tex lines 365--379, one \texttt{proof} environment of 15 lines). No mathematics is written, repaired or re-derived: the statement and the proof are the article's own lines, in the article's own notation, and no retained line is substituted or re-flowed.  No further source-local context is needed: every cross-reference of every retained line resolves inside the two delivered spans.  Nothing was omitted from any retained span, so the package carries no omission record.  No retained line contains a citation, so the wrapper carries no bibliography and no bibliography entry.  No retained line contains a figure environment, an image-inclusion command or drawing code, so no image is delivered and the package declares no asset dependency.  Editorial formatting only, in addition: an AMS \texttt{amsart} preamble that restates the article's own declarations and macros used by the retained lines, namely the environments \texttt{theorem} on separate counters, together with the standard packages those lines need.  The retained environments print in this document as Theorem 1 in order of appearance, whereas the article prints the same items with the numbering of its own full document.  The complete native source of the article as distributed in the arXiv e-print for this exact version is bundled unchanged as \texttt{sources/156375/source0.tex}.
\begin{theorem}
\label{pair}
There exists a pair of line arrangements in the complex projective plane, each of which has weak-combinatorics of the form
$$(d_{1};n_{2},n_{3},n_{4},n_{5},n_{6}) = (15; 24 ,12, 0, 0, 3),$$
where $n_{i}$ denotes the number of (ordinary) $i$-fold intersection points, such that one is free and the second is nearly-free.
\end{theorem}

\begin{proof}
The first line arrangement $\mathcal{L}_{1}$ is constructed by removing from the full monomial arrangement $\mathcal{FMA}_{6}$ defined by
$$F_{6}(x,y,z) = xyz(x^{6}-y^{6})(y^{6}-z^{6})(x^{6}-z^{6}) =0$$
the following lines
$$\ell_{1} : x-z=0, \quad \ell_{2} : x-ez=0, \quad \ell_{3} : y-z=0, $$
$$\ell_{4} : y-ez=0, \quad \ell_{1} : x-e^{2}y=0, \quad \ell_{6} : x-e^{4}y=0,$$
where $e^{2}-e+1=0$.
We can check, using \verb}SINGULAR}, that $\mathcal{L}_{1}$ is free with exponents $(d_{1},d_{2}) = (7,7)$. 

The second line arrangement $\mathcal{L}_{2}$ is constructed as follows. Take the full monomial arrangement $\mathcal{FMA}_{5}$ given by 
$$F_{5}(x,y,z) = xyz(x^{5}-y^{5})(y^{5}-z^{5})(x^{5}-z^{5}) =0$$
and remove the following three lines
$$\ell_{1} : x-z = 0, \quad \ell_{2} : x-y = 0, \quad \ell_{2} : y-z = 0.$$
We can check that $\mathcal{L}_{2}$ is not free, but only nearly-free with exponents $(d_{1},d_{2},d_{3}) = (6,9,9)$.
\end{proof}

\end{document}
